% !TeX root = ../main.tex
\section{Покрытие исходника и исправления}
Ниже указано, куда перенесён материал каждой страницы исходного PDF.
Повторы объединены, но ни одно самостоятельное решение не исключено.
Номер Т.1, Т.2 или Т.3 означает дополнительную задачу без номера
из задачника, помеченную в исходнике буквой «Т».
\begin{center}
\begin{tabular}{r|l}
Страница PDF & Задачи\\\hline
1 & 55.16, 55.21\\
2--3 & 55.29, 56.8, 56.18, а\\
4 & 56.18, а; 3.2; 3.4, б\\
5--6 & 3.7; 3.16, а; 56.3; 56.6; 56.10\\
7--8 & Т.1; 56.32, г; 56.37; 56.38\\
9--10 & 56.38; 56.40; 56.43; 56.44; Т.2, а\\
11--12 & 58.4, в; 58.5; 58.12; 58.16; 58.28, б\\
13--14 & 58.28, г; 58.29; 58.30; 58.31; 58.33; 55.25; 57.1\\
15--16 & 57.1; 57.3, а; 57.19; 58.7; 58.35; Т.3\\
17--18 & 60.4; 60.6; 56.41; 57.7\\
19--20 & 57.7; 57.14, б; 57.13, а, б; 57.38, б\\
21--22 & 57.38, б; 60.15; 60.17; 58.24, а\\
23--25 & 58.24, а, ж; 55.23; 55.32, б; 57.40; 57.45\\
26--27 & 56.18, б; 3.11; 3.16, б; 3.23; 56.28\\
28--30 & Т.2, б; 57.13, в; 58.36; 58.37; 58.21
\end{tabular}
\end{center}

\subsection*{Исправления, меняющие ответ или доказательство}
\begin{enumerate}
\item 56.8, б: доказано существование нужных степеней для любых
 конечных порядков, а не только для взаимно простых.
\item 3.11: исправлено направление изменения числа инверсий.
\item 3.16, а: исправлены формулы произведений тройных циклов
 при принятой композиции справа налево.
\item 3.23: ошибочное разделение пересекающихся транспозиций
 заменено явной конструкцией двух инволюций.
\item Т.1: к необходимости критерия добавлена достаточность
 и алгоритм построения квадратного корня.
\item 56.32, г: исправлена двойная транспозиция; отсутствие
 подгруппы порядка $6$ доказано отдельно.
\item 56.44 и Т.2: доказательства распространены на бесконечные
 группы конечного индекса без деления бесконечных мощностей.
\item 58.4, в: доказана полнота классификации нормальных подгрупп
 $S_4$ с помощью классов сопряжённости.
\item 58.12: заменена неверная классификация нормальных
 подгрупп диэдральной группы.
\item 58.30, б: степень $4$ заменена на $3$ для получения ядра $U_3$.
 В пункте г буквальный знаменатель $\Z_+$ не задаёт подгруппу;
 отдельно отмечена исправленная постановка с $\Rp$.
\item 58.31 и 58.33: записаны корректные проекции и
 гомоморфизмы вместо необоснованного «нормирования».
\item 55.25: для полной группы $\Z_3$ дана фигура без отражений.
\item 57.7: комплексные эрмитовы формы классифицируются по инерции,
 а не только по рангу; число орбит равно $\binom{n+2}{2}$.
\item 57.14, б: исправлена ось вращения икосаэдра.
\item 57.13, б: исправлено условие из задачника --- группа вращений
 куба равна $S_4$. В пункте в исправлен подсчёт $12\cdot5=60$
 и полностью обоснован изоморфизм с $A_5$.
\item Т.3: подсчитаны грани вместо рёбер; приведены ответы
 для вращений и для полной группы изометрий.
\item 60.6: аргумент трактуется по модулю $2\pi$.
\item 60.15, б: добавлено произведение с обратным элементом,
 необходимое для попадания в прямой сомножитель.
\item 58.24, ж: над $\R$ центр состоит только из $\lambda I$.
\item 58.21: исключено неправомерное предположение $ab=ba$.
\item 57.45: исправлены соотношения $D_4$ и степени автоморфизмов;
 вычислена вся подгруппа внутренних автоморфизмов.
\end{enumerate}

\subsection*{Источники и границы материала}
Основной источник --- все $30$ страниц «дз1\_чат\_кураторов.pdf».
Условия и терминология сверены с приложенным «Сборником задач
по алгебре» под редакцией А.~И.~Кострикина, издание 2009 года,
и лекционными материалами по теории групп.
Файл программы курса использован для проверки контекста заданий.
Задачи, отсутствующие в исходном PDF решений, не добавлялись
только на основании того, что они перечислены в программе курса.
Пункты, явно присутствующие в строках ответов или в напечатанном
условии рядом с решением, при необходимости кратко пояснены.

Оформление определения, теоремы, леммы, утверждения и доказательства
следует предоставленной преамбуле: используются команды
\verb|\Def|, \verb|\thm|, \verb|\lem|, \verb|\utv|;
формулировки теорем не нумеруются, поля страницы равны $15$ мм.
